% GENERATED FILE. Edit canonical lessons, then run npm run generate:books.
% canonical-source-hash: fnv1a64:ac7327a6777d3792
% canonical-lessons: UR-C287-angara, UR-C287-sulagna, UR-C287-sharara, UR-C287-bharakna, UR-C287-jalna

\chapter{Fire!}
\label{ch:ur-fire}
\hlchaptermodality{287}

\begin{chapteropening}
\textbf{By the end of this chapter:} \emph{I can say where a fire started and how it spread.}

Complete the last lesson of chapter 287: I can say where a fire started and how it spread.
\end{chapteropening}

\section[\texorpdfstring{\ur{انگارا} (angārā)}{انگارا (angārā)}]{\ur{انگارا} (angārā) — an ember; a live coal}

\label{lesson:UR-C287-angara}



\begin{quote}
Before the new one: say the Urdu for to crush, then the Urdu for to steady oneself.
\end{quote}

\subsection*{You'll want to know: \ur{انگارا}}
\textbf{\ur{انگارا}} — \emph{angārā} — \textquotedblleft{}an ember; a live coal\textquotedblright{}.

\begin{quote}
\ur{چولہے} \ur{میں} \ur{ایک} \ur{انگارا} \ur{تھا۔}
\end{quote}

\begin{quote}
\emph{chūlhe meṅ ek angārā thā.}
\end{quote}

\begin{quote}
There was an ember in the stove.
\end{quote}

Masculine.

\begin{cousinweb}
\textbf{\ur{انگارا}} (\emph{angārā}) is inherited from Sanskrit \emph{áṅgāra}, \textquotedblleft{}glowing charcoal\textquotedblright{} (CDIAL 125), where CDIAL lists the Hindi cousin \emph{ãgārā}.
\end{cousinweb}

\begin{grammarlens}[title={What you've built}]
One more word for telling what happened.
\end{grammarlens}

\subsection*{Your turn}
\begin{itemize}
\raggedright
  \item \emph{Say it:} \emph{angārā}
  \item \emph{Say it:} \emph{angārā}, once more
  \item \emph{Say it:} \emph{chūlhe meṅ ek angārā thā.} — There was an ember in the stove.
  \item \emph{Recall:} say the Urdu for to crush, then the Urdu for to steady oneself, then say \emph{angārā} again
\end{itemize}

\begin{tcolorbox}[breakable,colback=teal!4,colframe=teal!35!black,title={Before you move on}]
What does \ur{انگارا} mean? (\textquotedblleft{}An ember; a live coal\textquotedblright{}.) Say it once more.
\end{tcolorbox}
\section[\texorpdfstring{\ur{سلگنا} (sulagnā)}{سلگنا (sulagnā)}]{\ur{سلگنا} (sulagnā) — to smoulder}

\label{lesson:UR-C287-sulagna}



\begin{quote}
Before the new one: say the Urdu for to steady oneself, then the Urdu for an ember.
\end{quote}

\subsection*{You'll want to know: \ur{سلگنا}}
\textbf{\ur{سلگنا}} — \emph{sulagnā} — \textquotedblleft{}to smoulder\textquotedblright{}.

\begin{quote}
\ur{چولہے} \ur{میں} \ur{آگ} \ur{سلگ} \ur{رہی} \ur{تھی۔}
\end{quote}

\begin{quote}
\emph{chūlhe meṅ āg sulag rahī thī.}
\end{quote}

\begin{quote}
A fire was smouldering in the stove.
\end{quote}

Past (it …): \ur{وہ} \ur{آگ} \ur{سلگی}, \emph{voh āg sulgī}.

\begin{cousinweb}
\textbf{\ur{سلگنا}} (\emph{sulagnā}) is inherited: CDIAL lists the Hindi cousin \emph{sulagnā}, \textquotedblleft{}to be set on fire\textquotedblright{}, under Sanskrit \emph{saṃlagna}, \textquotedblleft{}attached\textquotedblright{} (CDIAL 12999). It is the slow burning before the flames come.
\end{cousinweb}

\begin{grammarlens}[title={What you've built}]
One more word for telling what happened.
\end{grammarlens}

\subsection*{Your turn}
\begin{itemize}
\raggedright
  \item \emph{Say it:} \emph{sulagnā}
  \item \emph{Say it:} \emph{sulagnā}, once more
  \item \emph{Say it:} \emph{chūlhe meṅ āg sulag rahī thī.} — A fire was smouldering in the stove.
  \item \emph{Recall:} say the Urdu for to steady oneself, then the Urdu for an ember, then say \emph{sulagnā} again
\end{itemize}

\begin{tcolorbox}[breakable,colback=teal!4,colframe=teal!35!black,title={Before you move on}]
What does \ur{سلگنا} mean? (\textquotedblleft{}To smoulder\textquotedblright{}.) Say it once more.
\end{tcolorbox}
\section[\texorpdfstring{\ur{شرارہ} (sharārā)}{شرارہ (sharārā)}]{\ur{شرارہ} (sharārā) — a spark}

\label{lesson:UR-C287-sharara}



\begin{quote}
Before the new one: say the Urdu for an ember, then the Urdu for to smoulder.
\end{quote}

\subsection*{You'll want to know: \ur{شرارہ}}
\textbf{\ur{شرارہ}} — \emph{sharārā} — \textquotedblleft{}a spark\textquotedblright{}.

\begin{quote}
\ur{چولہے} \ur{سے} \ur{ایک} \ur{شرارہ} \ur{نکلا۔}
\end{quote}

\begin{quote}
\emph{chūlhe se ek sharārā niklā.}
\end{quote}

\begin{quote}
A spark flew out of the stove.
\end{quote}

Masculine.

\begin{cousinweb}
is usually traced, through Persian, to Arabic \emph{sharāra}, \textquotedblleft{}a spark\textquotedblright{}. One \textbf{\ur{شرارہ}} is enough to start a fire.
\end{cousinweb}

\begin{grammarlens}[title={What you've built}]
One more word for telling what happened.
\end{grammarlens}

\subsection*{Your turn}
\begin{itemize}
\raggedright
  \item \emph{Say it:} \emph{sharārā}
  \item \emph{Say it:} \emph{sharārā}, once more
  \item \emph{Say it:} \emph{chūlhe se ek sharārā niklā.} — A spark flew out of the stove.
  \item \emph{Recall:} say the Urdu for an ember, then the Urdu for to smoulder, then say \emph{sharārā} again
\end{itemize}

\begin{tcolorbox}[breakable,colback=teal!4,colframe=teal!35!black,title={Before you move on}]
What does \ur{شرارہ} mean? (\textquotedblleft{}A spark\textquotedblright{}.) Say it once more.
\end{tcolorbox}
\section[\texorpdfstring{\ur{بھڑکنا} (bhaṛaknā)}{بھڑکنا (bhaṛaknā)}]{\ur{بھڑکنا} (bhaṛaknā) — to flare up}

\label{lesson:UR-C287-bharakna}



\begin{quote}
Before the new one: say the Urdu for to smoulder, then the Urdu for a spark.
\end{quote}

\subsection*{You'll want to know: \ur{بھڑکنا}}
\textbf{\ur{بھڑکنا}} — \emph{bhaṛaknā} — \textquotedblleft{}to flare up\textquotedblright{}.

\begin{quote}
\ur{ہوا} \ur{سے} \ur{آگ} \ur{بھڑک} \ur{اٹھی۔}
\end{quote}

\begin{quote}
\emph{hawā se āg bhaṛak uṭhī.}
\end{quote}

\begin{quote}
The wind made the fire flare up.
\end{quote}

Past (it …): \ur{وہ} \ur{آگ} \ur{بھڑکی}, \emph{voh āg bhaṛkī}.

\begin{cousinweb}
\textbf{\ur{بھڑکنا}} (\emph{bhaṛaknā}) is inherited: CDIAL lists the Hindi cousin \emph{bhaṛaknā}, \textquotedblleft{}to be excited, to take fright\textquotedblright{}, under \emph{*bhaṭ} (CDIAL 9365), with \emph{bhaṛak}, \textquotedblleft{}a flash\textquotedblright{}; a fire that flares is thought to share that sudden start. \textbf{\ur{بھڑک} \ur{اٹھی}} is the leap of the flames.
\end{cousinweb}

\begin{grammarlens}[title={What you've built}]
One more word for telling what happened.
\end{grammarlens}

\subsection*{Your turn}
\begin{itemize}
\raggedright
  \item \emph{Say it:} \emph{bhaṛaknā}
  \item \emph{Say it:} \emph{bhaṛaknā}, once more
  \item \emph{Say it:} \emph{hawā se āg bhaṛak uṭhī.} — The wind made the fire flare up.
  \item \emph{Recall:} say the Urdu for to smoulder, then the Urdu for a spark, then say \emph{bhaṛaknā} again
\end{itemize}

\begin{tcolorbox}[breakable,colback=teal!4,colframe=teal!35!black,title={Before you move on}]
What does \ur{بھڑکنا} mean? (\textquotedblleft{}To flare up\textquotedblright{}.) Say it once more.
\end{tcolorbox}
\section[\texorpdfstring{\ur{جلنا} (jalnā)}{جلنا (jalnā)}]{\ur{جلنا} (jalnā) — to burn; to be on fire}

\label{lesson:UR-C287-jalna}



\begin{quote}
Before the new one: say the Urdu for a spark, then the Urdu for to flare up.
\end{quote}

\subsection*{You'll want to know: \ur{جلنا}}
\textbf{\ur{جلنا}} — \emph{jalnā} — \textquotedblleft{}to burn; to be on fire\textquotedblright{}.

\begin{quote}
\ur{دکان} \ur{جل} \ur{گئی۔}
\end{quote}

\begin{quote}
\emph{dukān jal gaī.}
\end{quote}

\begin{quote}
The shop burned down.
\end{quote}

Past (it …): \ur{وہ} \ur{گھر} \ur{جل} \ur{گیا}, \emph{voh ghar jal gayā}.

\begin{cousinweb}
is inherited from Sanskrit \emph{jválati}, \textquotedblleft{}burns, blazes\textquotedblright{} (CDIAL 5306). It is the partner of \textbf{\ur{جلانا}}: you \textbf{\ur{جلانا}} a fire, and a thing \textbf{\ur{جلنا}} by itself.
\end{cousinweb}

\begin{grammarlens}[title={What you've built}]
One more word for telling what happened.
\end{grammarlens}

\subsection*{Your turn}
\begin{itemize}
\raggedright
  \item \emph{Say it:} \emph{jalnā}
  \item \emph{Say it:} \emph{jalnā}, once more
  \item \emph{Say it:} \emph{dukān jal gaī.} — The shop burned down.
  \item \emph{Recall:} say the Urdu for a spark, then the Urdu for to flare up, then say \emph{jalnā} again
\end{itemize}

\begin{tcolorbox}[breakable,colback=teal!4,colframe=teal!35!black,title={Before you move on}]
What does \ur{جلنا} mean? (\textquotedblleft{}To burn; to be on fire\textquotedblright{}.) Say it once more.
\end{tcolorbox}
